\section{Exclusion of isotropic extremizers}
\label{sec:isotropic}

We exclude the isotropic alternative in
Proposition~\ref{prop:isotropic-reduction}. The argument has two
geometric steps. Three sufficiently separated slopes would force
more point entropy than saturation allows, so one can retain a slope
population in a thin strip. That strip supplies an earlier anisotropic aperture;
a two-slope projection estimate in the remaining block then forces
a conditional time rate greater than one. Throughout both changes
of scale we retain uniform conditional probability bounds on particle
fibers, as well as the corresponding entropy inequalities.

\begin{proposition}[Isotropic exclusion]
\label{prop:isotropic-exclusion}
Suppose that $0<\gamma<1$ and that the classical universal growth rate satisfies
$\beta>\gamma$. There is no saturated isotropic system of duration and horizon
$1$ furnished by Proposition~\ref{prop:isotropic-reduction}, with homogeneous
conditional time profile $\Ent(C_u\mid T)=su$, $0\leq u\leq1$.
\end{proposition}

Here and below, a limiting assertion about homogeneous probabilities is used
on the trimmed laws of Lemma~\ref{lem:homogenization}. In particular, it supplies
uniform probability bounds on the retained joint population. An exceptional
set of small probability in an untrimmed conditional law is not treated as a
uniform probability bound. We give the finite estimates that implement this
distinction in the present argument.

\subsection{Finite incidence counts and regular increments}

In the next two elementary estimates $\rho$ is a mesh size, and all coordinates
belong to a fixed bounded set. Constants may depend on that set, but not on
$\rho$. A point cell in dimension $d+1$ has coordinates $(c,x)\in\RR\times
\RR^d$, and a slope bin has representative $v\in\RR^d$ and diameter $O(\rho)$.
The intercept map is
\[
 \pi_v(c,x)=x-cv.
\]

\begin{lemma}[Counting intercepts on retained fibers]
\label{lem:isotropic-intercepts}
Suppose an incidence population uses at most $N$ point cells. For every
retained exact particle and slope bin incident to it, suppose there are at
least $m$ distinct retained time children in the same base cell and slope bin.
Suppose further that these children lie on one line of true slope within
$O(\rho)$ of the bin representative. Then, for each slope bin, all its retained
incidences use at most $CN/m$ intercept cells of side $\rho$.
\end{lemma}

\begin{proof}
Fix the slope bin. A particle's intercept $x-cv$ varies by $O(\rho)$ along its
retained children. Its distinct time children give distinct point cells.
For every occupied intercept cell $b$, choose one particle witnessing an
incidence in $b$. The particle's $m$ point cells have intercepts in a fixed
bounded enlargement of $b$. A point cell can belong to the enlargements of
only $C$ intercept cells, since $v$ is fixed. Counting the pairs consisting of
an occupied intercept cell and one of these witnesses gives $m\#\{b\}\leq CN$.
This counts point cells, and imposes no bound on the number or entropy of exact
particles. The estimate concerns the population before replica labels are
fixed; its support upper bound continues to hold on every subsequent subset.
\end{proof}

We will use the following explicit form of the star calculation in
Theorem~\ref{thm:finite-conditioning}. Let $q$ be a normalized point--slope law
with actual marginals $\mu,\nu$, and suppose
\begin{equation}
 q(p,v)\leq K\mu(p)\nu(v),\qquad \mu(p)\leq M^{-1}.
 \label{eq:isotropic-pair-comparison}
\end{equation}
Sample $r$ slopes independently conditional on $p$. The resulting law $q_r$
satisfies both comparisons
\begin{align}
 q_r(p,v_1,\ldots,v_r)&\leq K^r\mu(p)\prod_{i=1}^r\nu(v_i),
 \label{eq:isotropic-star-comparison}\\
 q_r(p,v_1,\ldots,v_r)&\leq
 K^{r-1}q(p,v_j)\prod_{i\ne j}\nu(v_i)\quad(1\leq j\leq r).
 \label{eq:isotropic-one-pair-comparison}
\end{align}
These follow by multiplying $q(v_i\mid p)\leq K\nu(v_i)$, keeping one factor
unestimated for the second inequality. If $\lambda$ is the slope-tuple marginal,
then the tuples on which $\lambda/\prod_i\nu(v_i)<L^{-1}$ have
$\lambda$-mass at most $L^{-1}$, by summing against the product probability.
Outside that set,
\begin{equation}
 q_r(p\mid v_1,\ldots,v_r)\leq K^r L\mu(p).
 \label{eq:isotropic-pinning}
\end{equation}
Thus fixing a typical tuple leaves at least $M/(K^rL)$ point cells. An event
of independent slope probability $\eta$ has star probability at most $K^r\eta$.
All applications below first trim pairs, discard point fibers of inadequate
retention, and replace comparison marginals by the actual retained marginals
using Theorem~\ref{thm:finite-conditioning}. Consequently
\eqref{eq:isotropic-pair-comparison} holds for the normalized retained law.
The fixed labels will always be precisely these two or three slope replicas;
no cyclic incidence configuration is used.

\begin{lemma}[Two elementary projection inequalities]
\label{lem:isotropic-grid-projections}
Let $E$ be a set of occupied $\rho$-cells in a fixed bounded region.
\begin{enumerate}
\item In $\RR^3$, let $v_1,v_2,v_3\in\RR^2$ be bounded slopes and put
$D=\abs{\det((1,v_1),(1,v_2),(1,v_3))}>0$. If $\pi_{v_i}(E)$ uses at most
$N_i$ cells, with bounded enlargements permitted, then
\begin{equation}
 \abs E^2\leq C D^{-6}N_1N_2N_3.
 \label{eq:isotropic-lw}
\end{equation}
\item In $\RR^2$, for two bounded scalar slopes $v_1,v_2$ with
$\abs{v_1-v_2}\geq\rho^a$, if their intercepts use at most $N_1,N_2$ cells,
then
\begin{equation}
 \abs E\leq C\rho^{-2a}N_1N_2.
 \label{eq:isotropic-two-projections}
\end{equation}
\end{enumerate}
\end{lemma}

\begin{proof}
For the first assertion choose a witness point in each occupied cell, and
let $U$ have columns $(1,v_i)$. Replace each $U^{-1}z$ by its $\rho$-grid cell,
obtaining a finite grid set $Y$. Since $\norm U\leq C$, each cell of $Y$
accounts for at most $C$ original cells, so $\abs E\leq C\abs Y$.
The map $\pi_{v_i}U$ annihilates the $i$th coordinate. On the other two
coordinates it is an invertible matrix of determinant of absolute value $D$
and inverse norm at most $C/D$. Thus each original intercept cell accounts
for at most $CD^{-2}$ cells of the coordinate projection of $Y$ omitting
coordinate $i$. Rounding $U^{-1}z$ adds only $O(\rho)$ to its intercept,
because $\pi_{v_i}U$ has bounded norm.

For completeness, the discrete Loomis--Whitney inequality
\citep{LoomisWhitney1949} used here is
\[
 \abs Y^2\leq\abs{\pi_{12}Y}\abs{\pi_{13}Y}\abs{\pi_{23}Y}.
\]
To see it, write $f_{ij}$ for the indicators of the three projected sets.
The sum
\[
 \sum_{x,y,z}f_{12}(x,y)f_{13}(x,z)f_{23}(y,z)
\]
majorizes $\abs Y$.
Cauchy--Schwarz first in $z$ bounds this sum by
$\sum_{x,y}f_{12}(x,y)\sqrt{b_x c_y}$, where
$b_x=\sum_z f_{13}(x,z)$ and $c_y=\sum_z f_{23}(y,z)$.
Cauchy--Schwarz in $(x,y)$ bounds its square by
$\abs{\pi_{12}Y}(\sum_x b_x)(\sum_y c_y)$.
Combining the three projection bounds proves \eqref{eq:isotropic-lw}.
For the second assertion, the inverse of
$(c,x)\mapsto(x-v_1c,x-v_2c)$ has norm at most $C\rho^{-a}$.
Each pair of intercept cells therefore meets at most $C\rho^{-2a}$ point
cells. Summing over the pairs proves \eqref{eq:isotropic-two-projections}.
\end{proof}

Here is the scale-selection fact needed twice below. Let $F_u$, $0\leq u\leq1$,
be a filtration and let $X_h$ have normalized entropy at most $Dh$. For an
optional fixed conditioning variable $Z$, put
$g_h(u)=\Mut(X_h;F_u\mid Z)$. Nestedness and the chain rule give
\begin{equation}
 \int_0^{1-h}\Mut(X_h;F_{u+h}\mid Z,F_u)\,du
 =\int_0^{1-h}(g_h(u+h)-g_h(u))\,du
 \leq h\Ent(X_h\mid Z)\leq Dh^2.
 \label{eq:isotropic-integrated-information}
\end{equation}
The middle inequality follows by canceling the two integrals on their common
interval. For a sequence $h_j$ with $\sum_j\sqrt{h_j}<\infty$, Markov's
inequality and the elementary Borel--Cantelli argument show that, for almost
every $u$, the integrand divided by $h_j$ is at most $\sqrt{h_j}$ eventually.
The same points may be required to be differentiability points of any fixed
finite collection of Lipschitz entropy profiles.

Our anchored observations are nested only to bounded ambiguity. Formula
\eqref{eq:isotropic-integrated-information} is applied to their limiting
profiles. Its exact finite counterpart can instead use the nested dyadic
cells of the physical point $(c,A+cV)$: at every fixed depth $u$ these and
$(C_u,[A+C_uV]_u)$ determine each other with a bounded number of possibilities.
All entropy differences and conditional mutual informations consequently
differ by $O(1/\log(1/\eps))$. This error is taken to zero before $h$.

\subsection{The isotropic entropy bounds}

For a vector $Y$, write $[Y]_d$ or $Y_d$ for its isotropic grid
observation at scale $\eps^d$.
Assume the system in Proposition~\ref{prop:isotropic-exclusion} exists. Write
\[
 O_u=(C_u,[A+C_uV]_u),\qquad O=O_1,\qquad
 \alpha=2-2s+\beta.
\]
The exact particle $T$ determines $A,V$. In an anchored observation it leaves
only the time label random, so $\Ent(O_u\mid T)=su$. Write $\tau$ for the
root supremum, equivalently its uniformized cap level in the limiting system.
The terminal weight is zero. Saturation therefore gives
\begin{equation}
 \beta=-\Mut(T;O)+2s-\tau=3s-\Ent(O)-\tau.
 \label{eq:isotropic-terminal-score}
\end{equation}
Since $\tau\geq0$ and $\Mut(T;O)\geq0$, we have $\beta\leq2s$ and
\begin{equation}
 \Ent(O)\leq3s-\beta,\qquad
 \Ent(V_d\mid O)\geq\alpha d\quad(0\leq d\leq1).
 \label{eq:isotropic-alpha}
\end{equation}
To prove the second inequality, test at duration $1-d$ with velocity depths
$(d,d)$ and position depths $(1,1)$. Denote this observation by $\widetilde O$.
Its weight is $2d$, and its conditional time entropy is $s(1-d)$.
Universality and \eqref{eq:isotropic-terminal-score} imply
\[
 \Mut(T;\widetilde O)-\Mut(T;O)\geq(2-2s+\beta)d.
\]
The pair $(O,V_d)$ determines $\widetilde O$ to bounded ambiguity: transport
over a time displacement $O(\eps^{1-d})$ with velocity uncertainty
$O(\eps^d)$ has position error $O(\eps)$. The information difference is
therefore at most $\Ent(V_d\mid O)$ in the limit. At $d=1$ the identical
inequality follows from the definition of the root supremum at duration zero.
Notice that $\alpha\geq\beta>0$.

Homogenize jointly $O$, $V_d$, and their conditioning variables, for the
finite lists of depths used below. Formula~\eqref{eq:isotropic-alpha} then
means, on the trimmed joint population,
\begin{equation}
 \PP(V_d=q\mid O=o)\leq\eps^{\alpha d-\eta}
 \label{eq:isotropic-velocity-mass}
\end{equation}
with arbitrarily small $\eta>0$ after taking the inner limit. Indeed the
conditional probability exponent is the difference of the homogeneous
joint and marginal exponents, and that difference is at least $\alpha d$.
The same construction records the time probabilities conditional on the
exact particle. It is legitimate on saturated homogeneous parts by the
splitting and cap conclusions of Proposition~\ref{prop:isotropic-reduction}.

\subsection{Concentration of slopes in a thin strip}

The profile $u\mapsto\Ent(O_u)$ is nondecreasing and $3$-Lipschitz, starts at
zero, and has endpoint at most $3s-\beta$. There is consequently a set of
positive measure of regular points $k\in(0,1)$ with
\begin{equation}
 p:=\partial_k\Ent(O_k)<3s-\beta/2.
 \label{eq:isotropic-low-increment}
\end{equation}
Choose such a point satisfying \eqref{eq:isotropic-integrated-information}
with $X_h=V_h$, whose entropy is at most $2h$. Along a sequence $h\downarrow0$,
\begin{equation}
 \Ent(O_{k+h}\mid O_k)=(p+o(1))h,\qquad
 \Mut(V_h;O_{k+h}\mid O_k)=o(h).
 \label{eq:isotropic-first-increment}
\end{equation}

Inside a typical $O_k$-cell translate time and position to the cell center
and dilate both by $\eps^{-k}$. Put $\rho=\eps^h$, and call the child point
observation $P$. Bounded changes from anchored to physical point cells do
not affect any exponent. Homogeneity, the conditional fiber rule, and
Lemma~\ref{lem:zero-information}, followed by the actual-marginal conclusion
of Theorem~\ref{thm:finite-conditioning}, give a normalized retained pair law
$q(P,V_h)$ satisfying the following bounds, with $e=e(h)\to0$:
\begin{align}
 \#\supp P&\leq\rho^{-(p+e)},&
 \mu(P=z)&\leq\rho^{p-e},&
 q&\leq\rho^{-e}\mu\otimes\nu.
 \label{eq:isotropic-first-graph}
\end{align}
These conclusions hold outside base cells of probability tending to zero.
One may enlarge $e(h)$ to absorb all fixed numerical constants and all
$o(1)$ errors in \eqref{eq:isotropic-first-increment}. To be explicit about
the conditioning: the mean information in block units tends to zero; discard
base cells where it exceeds its square root, apply the likelihood-ratio
truncation in each remaining cell, and discard point fibers losing more than
half their mass. Homogeneous joint and marginal bounds then supply the two
point estimates in \eqref{eq:isotropic-first-graph}. This procedure introduces
only another quantity tending to zero into $e$.

We also trim on the exact particle fibers before sampling replicas. A retained
time child, conditional on $T,C_k$ in the original base law, has probability
at most $\rho^{s-e}$. Lift all pair restrictions to that law, and discard
particle fibers whose conditional retention there is below $\rho^e$.
If the pair restrictions have total retention $m_0$, these discarded fibers
have normalized retained mass at most $\rho^e/m_0$, by
\eqref{eq:retention-fibers}. Here $m_0$ is bounded below before this step;
more generally a subpower lower bound is absorbed by increasing $e$.
Each remaining particle has at least $\rho^{-s+2e}$ distinct retained time
children, by dividing its retained mass by the original child mass bound.
Given $T,C_k$, the anchored base observation is fixed. Thus this is a statement
in the very same base cell used for $P$; no particle entropy bound is involved.
For a slope bin of size $\rho$, all these points have intercept within $O(\rho)$
of the particle's intercept. Lemma~\ref{lem:isotropic-intercepts} gives
\begin{equation}
 \#\pi_v(\text{retained neighbors of }v)
 \leq C\rho^{-(p-s+3e)}.
 \label{eq:isotropic-first-intercepts}
\end{equation}
Apply one final simultaneous marginal core to this pruned law. Its retention
is at least one half, by Lemma~\ref{lem:marginal-core}, and the pair comparison
again uses its actual point and slope marginals. Increasing $e$ absorbs the
finite factors, so \eqref{eq:isotropic-first-graph} and
\eqref{eq:isotropic-first-intercepts} hold together for this final pair law.
The core need not preserve the rich particle fibers: the intercept support
bound was established before it, and survives by support inclusion. Use
this same order of operations in the two-replica application below.

Fix $\vartheta>0$ so small that $12\vartheta<\beta/8$. We claim that, in every retained
base cell, some strip of thickness $C\rho^\vartheta$ has slope mass at least
$\rho^{5e}$, after enlarging $e(h)\to0$ once more. If all such strips had
mass less than $\rho^{5e}$, two independent slopes would have distance at least
$\rho^\vartheta$ except with probability $C\rho^{5e}$; a ball of that radius is
contained in a strip of comparable thickness. Conditional on the first two,
the third would lie at distance at least $\rho^\vartheta$ from their line except
with the same probability. Hence, outside an independent event of probability
$C\rho^{5e}$,
\[
 \abs{\det((1,v_1),(1,v_2),(1,v_3))}\geq\rho^{2\vartheta}.
\]
For the three-replica star, \eqref{eq:isotropic-star-comparison} bounds the
exceptional probability by $C\rho^{2e}$. Choose $L=\rho^{-e}$ in
\eqref{eq:isotropic-pinning}. A good tuple then leaves a point population of
at least $\rho^{-p+5e}$ cells, and each projection has the support bound
\eqref{eq:isotropic-first-intercepts}. Lemma~\ref{lem:isotropic-grid-projections}
implies
\[
 2p-10e\leq3(p-s+3e)+12\vartheta+o_{\rho\to0}(1),
 \qquad\text{hence}\qquad
 p\geq3s-12\vartheta-19e-o_{\rho\to0}(1).
\]
This contradicts \eqref{eq:isotropic-low-increment} for all sufficiently small
$h$, with $\rho$ taken sufficiently small afterwards. We take $e>0$ even
when an error vanishes, so that the exceptional probabilities above tend to
zero in this inner limit.

Choose one such strip measurably in each base cell, using a finite grid menu
with bounded enlargement, and restrict incidences to it. The total retention
is at least $\rho^{5e}$ times the mass of the retained base cells. No strip
index is fixed globally: its index is a function of $O_k$ and will be included
in the earlier test through that base label. With
\begin{equation}
 l=\vartheta h,
 \qquad \abs{V\cdot n(O_k)-w(O_k)}\leq C\eps^l,
 \label{eq:isotropic-strip}
\end{equation}
round the unit normal and offset at precision $\eps^l$. This changes $C$ only.
The Radon--Nikodym bound for the restriction is $\eps^{-O(eh)}$.
Since $\vartheta$ is fixed and $e(h)\to0$, its normalized logarithmic cost is
$o(l)$.

The bins for $O$, the required $V_d$, and the particle-conditional time
probabilities are trimmed before this restriction. Lemma~\ref{lem:homogenization}
and its conditional retention rule
therefore preserve their exponents with errors $o(l)$, including the
particle-conditional time exponents and \eqref{eq:isotropic-velocity-mass}.
The absolute root cap bound gives $\tau_{\rm new}\leq\tau+o(l)$.
The terminal numerator before subtraction of $\tau$ changes by $o(l)$, because
its observation and conditional time exponents were recorded. Consequently
the new terminal score is at least $\beta-o(l)$, and universality bounds it
above by $\beta+o(l)$. We henceforth use this retained law and its own root
supremum; its terminal score is $\beta+o(l)$.

\subsection{An earlier anisotropic observation}

Take $h$ small enough that $k<1-l$, and put $a=1-l$. In the frame with normal
$n=n(O_k)$ and tangent $n^\perp$, form
\begin{equation}
 B=\bigl(O_k,O_a,[(A+C_aV)\cdot n]_1\bigr).
 \label{eq:isotropic-earlier-test}
\end{equation}
The coarse observation $O_a$ locates tangent position to depth $a$; its
normal position is additionally located to depth $1$. The strip supplies
normal velocity to depth $l$, and tangent velocity is bounded. The rounded
frame is known to depth $l$. Thus $B$, with boundedly many additional data
cells if necessary, is an admissible aperture with velocity depths $(0,l)$,
duration $a$, horizon $1$, and weight $(1-\gamma)l$.

It is determined by $O$ to bounded ambiguity. Indeed $O$ coarsens to $O_k,O_a$
to bounded ambiguity, hence determines the chosen strip up to the same fixed
number of choices. For its fine normal position use
\begin{equation}
 (A+C_aV)\cdot n=(A+C_1V)\cdot n+(C_a-C_1)w+O(\eps),
 \label{eq:isotropic-normal-transport}
\end{equation}
where the error is bounded by
$C\eps^a\eps^l=C\eps$. Conversely every component of the anchored $B$ is
fixed by $(T,C_a)$, since $C_a$ fixes $C_k$. In particular,
\begin{equation}
 \Ent(O\mid T,B)=sl+o(l).
 \label{eq:isotropic-last-time}
\end{equation}
The conditional time profile on the retained law follows from the ratios of
the recorded probabilities at depths $a$ and $1$; the strip restriction costs
$o(l)$ on typical particle fibers. Adjoining standard data cells to $B$
uses deterministic functions of $(T,C_a)$ and leaves this argument unchanged.

We prove the information saving
\begin{equation}
 \Mut(T;O)-\Mut(T;B)\geq sl-o(l).
 \label{eq:isotropic-information-saving}
\end{equation}
The main point is a lower bound for the last block's joint time and tangent
position entropy.

\subsection{Slope bounds inside the last block}

For $0\leq u\leq1$ let $F_u$ record, in addition to $B$, time to depth
$a+ul$ and tangent position at that time to the same depth. Within a $B$-cell
translate time and tangent position to their depth-$a$ centers and dilate
both by $\eps^{-a}$. The resulting coordinates are bounded, the new mesh
parameter is $\rho_*=\eps^l$, and the tangent slope
$X=V\cdot n^\perp$ remains bounded. The filtration $F_u$ has capacity $2$
per unit depth. Its base $F_0$ adds at most bounded information to $B$, and
$F_1$ is determined by $(O,B)$ to bounded ambiguity. Given $T,B$, its
conditional entropy profile is $su+o(1)$ in block units, because its new data
are fixed by the particle and the time cell.

For each fixed $d\in(0,1]$ and each fixed filtration depth $u$, the retained
conditional tangent slope law satisfies
\begin{equation}
 \PP\bigl(X\in J\mid B,F_u\bigr)
 \leq \rho_*^{\alpha d-o(1)}
 \quad\text{for intervals }J\text{ of length }\rho_*^d.
 \label{eq:isotropic-inherited-tangent}
\end{equation}
Here, as throughout, the inequality is uniform on a further trimmed
subprobability, followed by normalization on fibers of subpower retention.
We verify this transfer, since a bare conditional entropy inequality would
not suffice.

First apply \eqref{eq:isotropic-velocity-mass} at velocity depth $dl$ on the
jointly trimmed law. Adjoin the finite list of labels $(B,F_u)$ to $O$.
Each such list has at most $K$ possible values given $O$, where $K$ is fixed
when the list is fixed. For a positive tolerance $\zeta$, discard realized
values whose conditional likelihood given $O$ is below $\rho_*^\zeta$.
Their total probability is at most $K\rho_*^\zeta$. On the remaining values,
conditioning on the adjoined label multiplies a conditional velocity-bin
probability by at most $\rho_*^{-\zeta}$. More explicitly, write
$\kappa_h=O(e(h)/\vartheta)$ for an upper bound on all prior restriction costs
in last-block units. Whenever a retained subprobability has mass at least
$\rho_*^{\kappa_h}$, discard conditioning fibers whose retention is below
$\rho_*^{\kappa_h+\zeta}$. Formula~\eqref{eq:retention-fibers} bounds their
normalized retained mass by $\rho_*^\zeta$; normalization costs at most
$\kappa_h+\zeta$ in the conditional exponent. For the fixed finite number
of such operations, the total extra loss is $C\kappa_h+C\zeta$.
It tends to zero by first taking the inner limit, then $h\downarrow0$ and
$\zeta\downarrow0$. It incurs no factor involving the number of $O$-cells
or $B$-cells.

For a fixed $B$-cell, its normal strip has thickness $O(\eps^l)$. Its
intersection with a tangent interval of length $\eps^{dl}$ meets only $C$
velocity grid cells of depth $dl$, since $d\leq1$. Thus the same uniform
upper bound holds for this tangent interval conditional on $O,B,F_u$.
Averaging over compatible $O$-labels gives
\eqref{eq:isotropic-inherited-tangent}, with the explicit extra exponent
loss $\zeta$ and the retention losses. Taking those losses to zero proves
the asserted form. Every bound is imposed on the retained joint population
before averaging; no exponentially small bound for untrimmed exceptional
probabilities is needed.

We now make the order of limits precise. Fix $\vartheta$ first. Choose the regular
sequence $h\downarrow0$ above. For each fixed $h$ record the finite lists of
$O,V_d$, and particle-conditional time bins before the strip choice.
Construct $B,F_u$ and perform the transfer just proved after choosing the
strip. The later increment binnings are performed after this transfer.
Take $\eps$ small enough that all inner mesh and bin errors, divided by $l$,
tend to zero. A diagonal,
with $h\downarrow0$, gives last-block systems with parameter $\rho_*$, exact
limiting time profile $su$, and \eqref{eq:isotropic-inherited-tangent} at a
countable dense list of $(u,d)$. Bounded coordinate capacities extend the
entropy profiles to other depths. For the conditional mass bound at a
subsequently chosen fixed filtration depth $u$, repeat the direct
$O\to(B,F_u)$ determination and likelihood truncation above; it has the
same bounded multiplicity. A requested slope depth $d$ is obtained by
using recorded depths $d'<d$ and then letting $d'\uparrow d$.
Thus the mass bound is inherited from that comparison, independently
of continuity of entropy. A later increment length $r>0$ is held fixed before
each inner limit; only then is $r\downarrow0$. Thus errors $o(l)$ from the
first construction are never divided by an unchosen infinitesimal $r$.
No bound on $\Ent(B)/l$ is asserted or needed.

\subsection{Two-slope expansion and the score contradiction}

We show, in the last-block limits just constructed, that
\begin{equation}
 \Ent(F_1\mid B)\geq2s
 \quad\text{in units }\log(1/\rho_*).
 \label{eq:isotropic-last-entropy}
\end{equation}
Write $f(u)=\Ent(F_u\mid B)$ in these units. It is nondecreasing and
$2$-Lipschitz, with $f(0)=0$. Apply
\eqref{eq:isotropic-integrated-information} to $X_r$, the tangent slope
prefix of depth $r$, conditional on $B$. Its entropy is at most $r$, so the
integrated mutual information is at most $r^2$. At almost every regular
interior $u$ there is consequently a sequence $r\downarrow0$ for which the
increment has negligible information about $X_r$, in units of its own length.

In a typical $(B,F_u)$-cell let $P_1$ be the time--tangent-position increment,
at mesh $\rho=\rho_*^r$. Homogenize its probabilities and the pair with $X_r$.
Let $p_1$ denote its point exponent. The likelihood-ratio and actual-marginal
steps used in \eqref{eq:isotropic-first-graph} give, with $e\to0$ along the
regular sequence,
\[
 \#\supp P_1\leq\rho^{-(p_1+e)},\qquad
 \mu_1(z)\leq\rho^{p_1-e},\qquad
 q_1\leq\rho^{-e}\mu_1\otimes\nu_1.
\]
The particle-conditional time exponent is $s$. Trimming exact particle
fibers and applying Lemma~\ref{lem:isotropic-intercepts} therefore bounds
every retained intercept support by $C\rho^{-(p_1-s+3e)}$.

Fix $0<a_0<1$. Formula~\eqref{eq:isotropic-inherited-tangent}, at depth
$a_0r$, gives the actual retained slope marginal the interval bound
$\nu_1(J)\leq\rho^{\alpha a_0-o(1)}$ for intervals of length $\rho^{a_0}$.
The two-replica star consequently has
$\abs{v_1-v_2}\geq\rho^{a_0}$ outside a set of probability
$O(\rho^{\alpha a_0-2e-o(1)})$. Since $\alpha>0$, this probability tends
to zero by taking $r$ sufficiently far along the regular sequence and then
the inner small-parameter limit. Pin a good tuple using
\eqref{eq:isotropic-pinning} with $L=\rho^{-e}$. At least
$\rho^{-p_1+4e}$ point cells remain. Their original intercept support bounds
are still valid. Formula~\eqref{eq:isotropic-two-projections} now gives
\[
 p_1-4e\leq2(p_1-s+3e)+2a_0+o(1),
 \qquad p_1\geq2s-2a_0-10e-o(1).
\]
First let the regular-increment errors vanish, and then let $a_0\downarrow0$.
This proves the lower entropy rate $p_1\geq2s$.

Here is why the cellwise argument yields $f'(u)\geq2s$ almost everywhere,
even if new homogeneous parts have different point exponents. Their
conditional increment entropies average to $f(u+r)-f(u)$, and all have
capacity at most $2r+o(r)$. If $f'(u)<2s$ by a fixed amount, a positive
fraction of the conditional population must have point exponents below
$2s$ by a smaller fixed amount. The mean conditional information is $o(r)$,
so Markov's inequality removes a vanishing fraction while imposing the pair
comparison above. The uniform time and slope bounds survive on the remaining
typical fibers. Homogenizing a finite list costs a subpower, and thus selects
a part with all these properties and the same fixed deficit. The preceding
statement about inherited bounds can also be checked directly: in a fixed
conditional cell, restriction to a part $E$ of mass $m$ gives
$\nu_E(J)\leq m^{-1}\nu(J)$. Here $m=\rho^{o(1)}$, so the slope exponent
is preserved. Discard exact particle fibers retaining less than
$m\rho^\zeta$ of their original conditional time law; their normalized
retained mass is at most $\rho^\zeta$. The time exponent therefore loses
at most $\log_{1/\rho}(1/m)+\zeta=o(1)+\zeta$. These estimates justify
the simultaneous selection without any count of conditioning cells.
The two-slope inequality contradicts the fixed deficit. Since $f$ is Lipschitz,
integrating $f'(u)\geq2s$ proves \eqref{eq:isotropic-last-entropy}.

Returning to the original units, \eqref{eq:isotropic-last-entropy} and the
determination of $F_1$ by $(O,B)$ give
\[
 \Ent(O\mid B)\geq2sl-o(l).
\]
Together with \eqref{eq:isotropic-last-time}, this gives
$\Mut(T;O\mid B)\geq sl-o(l)$. Since $\Ent(B\mid O)=o(l)$,
the information chain rule proves \eqref{eq:isotropic-information-saving}.

Finally let $Q_G$ denote the undivided score numerator, including subtraction
of the root supremum of the retained law. The same root supremum occurs for
$B$ and $O$. Their time entropies differ by $sl+o(l)$, so
\begin{align*}
 Q_B-Q_O
 &=(1-\gamma)l+\Mut(T;O)-\Mut(T;B)-2sl+o(l)\\
 &\geq(1-\gamma-s)l-o(l).
\end{align*}
On the other hand terminal saturation and the universal duration bound give
$Q_O=\beta+o(l)$ and $Q_B\leq\beta(1-l)+o(l)$. Divide by $l$ and pass through
the specified limits to obtain
\[
 1-\gamma-s\leq-\beta,
 \qquad s\geq1+\beta-\gamma>1.
\]
This contradicts the scalar time capacity $s\leq1$ and proves
Proposition~\ref{prop:isotropic-exclusion}.
